\section*{Revision basis and scientific changes}
We thank the referee for distinguishing genuine scientific revisions from branch names and publication plans. The latest R56 report at review commit \path{1bf0400f3540605d2ac68d2a40e8bcfb377d5446} correctly identified the missing R56 manuscript. The last complete predecessor was R54 at \path{eee0d39e08eaedc3a26cad792b71c9c5b5a8abd9}.

R58 starts from the subsequent, readable R57 source commit \path{debd4dfa4053cf5f3d6fcf06633b749ef6144ed6}, which itself starts directly from the R54 scientific tip. R57's workflow executed its declared study and built its readers, but failed while writing archive metadata, before the final repository publication. We do not describe that workflow as a completed publication. R58 repairs the packaging and publication checks, supplies additional structural results, and executes a new source-frozen study. Its numerical statements are generated from the new records, not copied from the failed run.

The new structural results are an oscillation-certified approximate-type compression theorem and an exact all-budget saturated-promise frontier. They strengthen the small-menu and heterogeneous-reward answer while preserving the recovered cap-count, deficit, lattice and price-path work. Reviewer-output commits are not added to the author-revision delta. All inherited scientific-source blobs remain unchanged outside the explicitly versioned root readers, whose exact predecessors are retained. Existing branches remain unchanged.

The original optimization problem is retained: accepted promise, expected caps, hard realization ceilings, pre-draw service, common paid catalog, nonnegative opening charges and command budget. No representative-history relaxation, enlarged budget, rounded obligation or deletion of adverse historical evidence is introduced.

\section{Latest R56 report: completeness and branch identity}
\subsection{A genuine reader, not a protocol-only branch}
The root \path{main.tex} assembles the revised title, text-only abstract, introduction, model, constructive representations, complexity results, algorithmic guarantees, computational discussion and references. Its mathematical dependency list is recorded by the TeX recorder, rather than inferred from directory names. The root \path{main.pdf} and \path{electronic_companion.pdf} are compiled from these exact entry points. This response is \path{RESPONSE_TO_REFEREES.pdf} in the R58 directory. The current README gives one scientific parent, one current revision and explicit reproduction commands.

The build validation checks the required inherited and new theorem labels, unresolved citations and cross-references, duplicate labels, overfull boxes, blank pages, source inputs and PDF hashes. The format is eleven point, one-and-one-half spacing and one-inch margins; the measured nonreference page count determines the journal's regular or lengthy category. Build success is a packaging check, not an editorial decision.

\subsection{A stable, inspectable claim set}
The original selected-boundary, packing, exact price-path, weak hardness, common-cap and safe-screening results are retained. The new cap-count theorem and its strict examples appear in Section~\ref{sec:smallmenus56}. Theorem~\ref{thm:types57} extends the constructive compression to joint cap--reward classes. The deficit representation, centered approximation, support-hole convolution and explicit cost bound are Theorems~\ref{thm:deficit56} and~\ref{thm:approx56}, Lemma~\ref{lem:holes56} and Proposition~\ref{prop:cost56}. Theorem~\ref{thm:lattice56} gives the exact zero-service lattice regime. These are statements with proofs and reconstruction, not protocol descriptions.

\subsection{Official execution and source binding}
R58 declares 186 primary timed method runs and 39 inherited-design extension runs. All are executed under the corrected executor; no favorable subset of earlier local timings is substituted. Inputs are serialized and hashed before their method calls. The source manifest is written before timed execution and checked again afterward. Raw run records, last phases, failures, rational certificate bytes, initial verification outcomes and separate extended checks are retained. Table-generation code derives numerical claims from these records. The publication audit refuses to certify completeness unless all 225 declared records exist and every published rational certificate has a successful independent check of its unchanged bytes.

The design was inspected during earlier local development. The current execution is therefore a reproducibility study on a disclosed design, not an unseen confirmatory experiment. Mechanical replay is separately identified and does not replace benchmark time. The missing, different R55 development archive is not fabricated from its protocol or counted as an R58 result. The original R55 protocol files remain on their original branch; the recovered local R56 package is identified by content hash in the R58 protocol. Claims in the current article rely on the actual R58 repository records.

\section{R54 and R56: small-menu structure and encoding complexity}
\subsection{A constructive and tight cap-count boundary}
Theorem~\ref{thm:capcount56} proves that, under a common reward, twice the number of distinct expected caps suffices. It allows heterogeneous realization ceilings and convex nondecreasing service costs. Its construction conditions a supplied original policy on each cap class, retains the class's weighted promise, and applies the common-cap singleton/pair reduction to that class. Taking the union preserves the original aggregate promise and all history-level constraints; nonnegative opening charges cannot rise after deletion. Proposition~\ref{prop:tight56} exhibits instances requiring all twice-as-many commands. Unused zero-charge commands need not be counted, and the theorem does not promise useful compression when almost every cap differs.

R58 adds Theorem~\ref{thm:types57}. Grouping jointly by expected cap and catalog reward type permits different rewards across groups. The same conditional construction then gives a bound of twice the number of joint groups. This is implemented for heterogeneous quadratic rewards and independently checked on 36 additional models with within-type heterogeneous costs and ceilings. It strengthens the structural answer rather than imposing a common-reward restriction on the heterogeneous model.

The construction uses a supplied feasible policy; it does not assume knowledge of the unknown optimum. To find an optimum without such knowledge, the paper separately gives exact enumeration up to the proven cardinality bound. Its exponent depends on the number of types or the command budget. We explicitly identify this as XP, not fixed-parameter tractability. No unproved FPT or parameterized-hardness assertion is used.

\subsection{An exact lattice algorithm at arbitrary command budgets}
Theorem~\ref{thm:lattice56} gives the requested pseudo-polynomial exact regime, including the zero-service subclass used by the weak SUBSET SUM reduction. Weighted command quantities, weighted caps and the promise must share a stated rational lattice. On a fixed selected path the zero-service kernels are concave piecewise linear with lattice segment lengths. Filling marginal segments produces an optimal allocation whose edge deficits lie on the same lattice. The merged recurrence then optimizes over all selected paths at the original command budget, with exact original-policy reconstruction.

The state parameter is the numerical lattice deficit, not its bit length. Under uniform history probabilities and a common command/cap grid, the required scale is explicit. General rational denominators can have an exponentially large numerical least common multiple. Consequently the algorithm is pseudo-polynomial in the scaled integral encoding, fully compatible with the encoding-sensitive weak hardness theorem. Positive quadratic service costs can introduce off-lattice stationary allocations; they remain covered by the additive theorem and exact branching, not by an unsupported lattice-exactness claim.

\section{R54 and R56: heterogeneous resource computation}
\subsection{A source-independent resource state}
The selected-boundary capacities are differences of one vertex potential. Their sum on an anchored book equals aggregate capacity minus the anchor. Subtracting used capacity from available capacity therefore makes total deficit equal aggregate capacity minus the accepted promise, independently of the anchor. Theorem~\ref{thm:deficit56} makes every feasible anchor an initial state at deficit zero and every completed book a state with the same terminal deficit. A predecessor still records the actual anchor for original-policy recovery. This removes an entire catalog factor from the earlier separately anchored recurrence.

The centered approximation uses the midpoint of the global marginal interval. Rounding deficits downward yields an upper certificate; restoring the original total deficit gives a feasible lower policy. Theorem~\ref{thm:approx56} bounds the interval by the requested additive tolerance without modifying any feasibility primitive. The history-level packing proof permits fully heterogeneous increasing concave rewards, while the implemented finite-arithmetic bound is stated for rational quadratic primitives.

\subsection{Support holes, bit work, memory and checker costs}
The recurrence does not assume that merged Bellman rows are concave or have connected reachable support. Lemma~\ref{lem:holes56} proves monotonicity of the smallest maximizing predecessor in the feasible staircase band, even after unreachable columns and rows are removed. Two hundred comparisons against exhaustive convolution test those support holes.

Proposition~\ref{prop:cost56} counts kernel preparation, selected counts, catalog edges and resource positions. It includes the edge cache and predecessor storage. Rational kernel evaluations and finite path sums have explicitly controlled binary encoding lengths. The cost is polynomial in input bit length and numerical grid size, not in the logarithm of inverse error or of the common lattice scale.

The deficit checker reconstructs kernels and verifies all Bellman upper inequalities through direct predecessor enumeration. It imports neither the deficit optimizer nor its fast convolution. Its conservative cost can be larger than the optimizer's cost; serialization and checking are measured separately. Price certificates use a backward at-most-budget recurrence independent of the optimizer's forward exact-count recurrence, together with explicit original lotteries and a complete leaf cover. The full large-integer regression, enormous-state pre-allocation rejection and seven mutation classes distinguish mathematical validity, arithmetic cost and encoding robustness.

\section{R54 and R56: exact search and end-to-end evidence}
\subsection{What the node oracle does and does not establish}
The contribution of exact price paths is the model-specific bound and its original-space certificate. Finite inclusion/exclusion branching and Lagrangian fixing remain credited to established optimization methods. The worst-case binary tree remains exponential; no runtime observation is used to claim otherwise.

The new records expose evaluated and created nodes, completed oracle calls, root and final gaps, depth, live leaves, incumbent updates, bound and infeasibility pruning, required-set sizes, fixed-book counts and time-to-bound trajectories. A node and a price call are not interchangeable units. Direct-tree forced fixings are distinguished from the separate screening procedure. The exact-target panel compares price search and enumeration at the identical zero-width request; the positive-tolerance and numerical-envelope panels retain their distinct completion definitions.

Nine untimed mechanical replays record tested prices and required/forbidden sets. The equal-book-count enumeration diagnostic measures incumbent selection, not equal arithmetic work. Retained-leaf price bit lengths and full raw/gzip proof bytes are reported. All failed historical R54 reduction challenges and large-history limits remain unchanged, alongside the new execution's unsuccessful cases.

\subsection{Screening under the same total allowance}
The 18 paired cases give direct price solving and screening followed by reduced solving the same total optimization allowance. Screening pays for its incumbent construction and forced-command bounds before spending any remaining time on the reduced problem. The report separates possible anchors, wholly ineligible commands and ordinary nonanchors, and records original-model interval width, serialization, verification and total proof time. Both favorable and unfavorable pairs are retained. The number of total-time improvements and paired ratios are generated from actual R58 records, not copied from earlier local outcomes. Correct exclusion and removal percentages are not asserted to imply a speed improvement.

\section{Operations Research contribution and preservation}
\subsection{Structural significance and operational interpretation}
The revision takes the structural route explicitly allowed by Section 5.8 of the latest report. It supplies a tight cap-count boundary, a finite heterogeneous-type extension, a source-independent resource representation, a faster heterogeneous approximation and an exact encoding-sensitive lattice regime. These are new mathematical and algorithmic statements on the original finite-catalog problem, not formatting or audit claims.

The companion also retains the service-credit implementation interface, contractual versus estimated inputs, pre-draw timing and charge interpretation. The payoff-calibration transfer result applies with feasibility inputs fixed; it is not misrepresented as robustness to changed probabilities or caps. The compatible-production example and its original evidence remain with the synthetic meaning of module return coefficients stated. No real-world field calibration or universal solver superiority is inferred from synthetic data.

\subsection{No arbitrary deletion or silent revision of old results}
The article retains the selected-boundary construction, heterogeneous component packing, weak hardness, exact priced-path oracle, common-cap tractability, additive predecessor, safe catalog rules and prior computational interpretation. Complete historical tables and supporting proofs remain in the companion or unchanged historical sources. The content map identifies where the new arguments enter and where antecedents remain. A full Git-blob preservation manifest checks every inherited file, with exact copies for replaced root readers. The code/data package includes runnable dependencies, inputs, new records and certificates; the Git tree retains older complete readers.

The remaining mathematical distinctions are local and explicit: the tight cap-only result requires a common reward, the joint-type extension groups by both cap and reward, the exact lattice theorem requires zero service costs, and exact branching may be exponential. These are not substitutes for revision. They accompany the new proofs, implemented algorithms, complete declared execution and independently checked original-policy certificates supplied for substantive re-review.

\section{Additional R58 structural answers and verification}
\subsection{Reward types need not coincide exactly}
Theorem~\ref{thm:robusttypes58} proves a menu-size versus approximation tradeoff in the original feasible set. It compares expectations of the catalog reward error under two feasible policies, so the loss is bounded by its weighted oscillation rather than an unnecessarily larger absolute-level error. Additive reward offsets cancel. The result supplies both a compression of any given policy and an optimization procedure through a prototype problem. No unknown optimum is required as algorithmic input. The implementation checks 24 original/prototype pairs against exhaustive enumeration, including zero-dispersion cases, and independently verifies the recovered original policies.

\subsection{A genuinely polynomial arbitrary-budget regime with positive service costs}
Corollary~\ref{cor:saturation58} exploits a saturated accepted promise, not zero service cost. The original target constraints force every target to its own cap. Every deficit is zero, so evaluating full-capacity arc weights and one selected-count recursion produces the complete at-most-budget value frontier. The method retains heterogeneous caps, ceilings, rewards and positive costs. A separate implementation is compared with exhaustive enumeration on 24 models across all seven budgets. All 168 original-space budget certificates are checked by the existing independent deficit verifier, which does not import the new frontier solver. Off-saturation requests and altered upper bounds are rejected.

\subsection{Publication failure repaired without altering the evidence}
The R57 failure called a file-size property as a function after the PDFs had already been built. The old archive also gathered unrelated historical evidence into one unnecessarily large ZIP. R58 supplies a checked, self-contained current-study and reader package with all original records used in its retained R52--R54 tables. Earlier repository evidence is preserved in Git, not deleted to make the archive smaller. Archive integrity, reader dependencies, current source hashes, the declared execution set and inherited Git blobs are checked before publication. A clean extracted copy rechecks the structural tests and rebuilds the readers. Automated PASS denotes these finite checks and artifact integrity; the universal proofs remain mathematical arguments submitted for referee review.
